\question[3]
Erläutere den Unterschied zwischen symmetrischer und asymmetrischer Verschlüsselung.  

\begin{solutionbox}{5cm}

Bei der symmetrischen Verschlüsselung wird für das Ver- und Entschlüsseln nur ein Schlüssel verwendet.
Der Schlüssel muss entweder persönlich oder mit Hilfe eines 
abhörsicheren Kanals vom Sender zum Empfänger übertragen werden.
Bei der asymmetrischen Verschlüsselung gibt es zwei Schlüssel, einen privaten und einen öffentlichen
Schlüssel, der letztere ist allen möglichen Kommunikationspartnern bekannt. Der private Schlüssel wird
nie veröffentlicht, der öffentliche Schlüssel kann über alle zur Verfügung stehenden Publikationsmöglichkeiten
veröffentlicht werden (E-Mail, WWW-Seiten, Public-Key-Server, …).
Um die Echtheit eines öffentlichen Schlüssels zu überprüfen, muss es ein Verfahren zur Herstellung von
Schlüsselvertrauen geben. Das kann eine (zentrale) Zertifizierungsstelle sein oder ein dezentraler
Web-of-Trust-Ansatz.

\end{solutionbox}
